\begin{lemma}\label{lem:nonic117-level39-fibres}
Over $\mathbb F_{512}$, every elliptic curve with $468$ rational points,
together with a point of order $39$, gives one of the four known
closed nonic points of $X_1(39)$.
\end{lemma}
\begin{proof}
The complete ordinary elliptic-curve census gives exactly three
$j$-invariants with a twist of order $468$.
In the common field $\mathbb F_2[t]/(t^9+t^8+1)$ their integer
coefficient encodings are $130,315,441$.
Each has precisely one twist of this order.
Since $\gcd(468,511)=1$, each rational point group is cyclic, and
there are $\varphi(39)/2=12$ fine marked structures per curve.
Their $j$-invariants are nonzero, so the only automorphisms are
$\pm1$, which act freely on the points of order $39$.

The four known nonic points have disjoint Frobenius orbits of length
$9$. Their $36$ exact Tate pairs give precisely $12$ pairs at each
of these three $j$-invariants. The binding certificate checks this
directly from the normalized pairs, and independently counts $468$
points on each representative Tate curve. Thus these points exhaust
all marked structures under consideration.
There is no supersingular exception: a supersingular elliptic curve
in characteristic $2$ has no geometric point of order $2$, so its
finite group has odd order and cannot have order $468$.
\end{proof}
